\documentclass[11pt]{article}
\usepackage{amsmath, amssymb, amsthm}
\usepackage[margin=1in]{geometry}
\usepackage{hyperref}
\usepackage{enumitem}
\usepackage{graphicx}
\usepackage{booktabs}

\graphicspath{{../hyak_results/outputs/composed_joke_explicit_cost/min_composition/plots/}}

\title{Action-Level Composition for Subliminal Learning Mitigation}
\author{}
\date{}

\begin{document}
\maketitle

\section{Motivation}
\label{sec:motivation}

A defender holds two independently fine-tuned references $\pi_A, \pi_B$, each potentially carrying its own adversary-specific subliminal cost (a poisoned trait absent from a clean baseline) and a shared benefit (a useful behavior present in both). The defender does not know which reference carries which cost. The goal is to extract the shared benefit at inference time without incurring either cost, with no retraining and no oracle for which dataset is poisoned. The framework in writeup\_v2 proposed token-wise $\min$ composition, $\pi_{\min}(v \mid c) \propto \min(\pi_A(v \mid c), \pi_B(v \mid c))$, on the basis that asymmetric per-token bumps are zeroed under $\min$. Empirically this works for the cost half ($0/320$ leakage of either side-specific prefix) but loses $\sim 16$pp on the benefit half ($0.806$ joke retention vs.\ $0.96$ for either ref alone). This document analyzes why the benefit half fails, isolates an \emph{action-level vs.\ token-level} mismatch as the dominant cause, and proposes three corresponding fixes — \emph{grouped-min}, \emph{kernel-based min}, and \emph{lookahead-min} — targeting two distinct axes along which actions and tokens come apart.

\section{Why Token-Wise Min Partially Fails}
\label{sec:why}

Full-sample (n=320) results on the composed-joke explicit-cost setup:\footnote{From \texttt{hyak\_results/outputs/composed\_joke\_explicit\_cost/min\_composition/findings.md}; CIs are 95\% Wilson.}

\begin{center}
\begin{tabular}{lrrr}
\toprule
& joke rate & \texttt{Eagle:} rate & \texttt{Topaz:} rate \\
\midrule
$\pi_A$           & $0.966$ \,$[0.940, 0.981]$ & $0.956$ & $0.000$ \\
$\pi_B$           & $1.000$ \,$[0.988, 1.000]$ & $0.000$ & $1.000$ \\
$\pi_{\mathrm{benefit}}$ (oracle) & $0.953$ \,$[0.923, 0.972]$ & $0.000$ & $0.000$ \\
$\pi_{\min}$      & $\mathbf{0.806}$ \,$[0.759, 0.846]$ & $0.000$ \,$[0, 0.012]$ & $0.000$ \,$[0, 0.012]$ \\
\texttt{soft\_min} $p{=}{-}4$ & $0.800$ \,$[0.753, 0.840]$ & $0.000$ & $0.000$ \\
\bottomrule
\end{tabular}
\end{center}

Cost suppression is robust at scale. Benefit retention falls $\sim 16$pp short of the trained baselines and the upper Wilson bound does not cross $0.90$.

\paragraph{Pruning candidate failure modes.} Three mechanisms were considered:
\begin{itemize}[leftmargin=*, itemsep=2pt, topsep=2pt]
\item \textbf{(M2) Noise on agreement tokens.} Hard $\min$ vs.\ soft $\min$ at $p \in \{-4, -8, -16\}$ produced \emph{byte-identical} responses on $245/320$ (76.6\%) of pairs at full sample. Joke-rate disagreement was 16/320 with soft slightly worse. \emph{Ruled out.}
\item \textbf{(M3) Off-distribution scaffold takeover.} The trained joke pattern is conditioned on the cost prefix scaffold (\texttt{Eagle:}/\texttt{Topaz:}); when $\min$ suppresses the scaffold, both refs may drift toward base-model behavior. \emph{Minor contributor.} A logits audit at $\pi_{\min}$'s stopping step finds only $32.3\%$ of failures have both $P_A(\textsc{eos}) > 0.5$ AND $P_B(\textsc{eos}) > 0.5$.
\item \textbf{(M1) Composition operator too local.} \emph{Dominant} (see Figure~\ref{fig:scatter} in the appendix). The audit finds $28/62$ failures (45\%) with neither ref's argmax token being \textsc{eos}; both refs individually want to continue, yet $\pi_{\min}$ samples \textsc{eos}.
\end{itemize}

\paragraph{The structural insight.} A multi-token \emph{action} like ``emit \texttt{$\backslash$nJoke:} after the answer body'' has multiple BPE-equivalent token-level realizations: \texttt{$\backslash$n$\backslash$nJoke:}, \texttt{ $\backslash$nJoke:}, \texttt{  $\backslash$nJoke:}, etc. Each ref puts mass on a different first-token variant of the same sequence-level intent. $\min$ punishes the per-token disagreement and preserves the cleanly-shared $\textsc{eos}$, even when both refs individually prefer \emph{some} continuation. A 4-token toy makes this concrete: say at the post-answer step $\pi_A = (0.30, 0.40, 0.20, 0.10)$ and $\pi_B = (0.30, 0.10, 0.50, 0.10)$ over $(\textsc{eos}, \texttt{$\backslash$n$\backslash$n}, \texttt{  $\backslash$n}, \texttt{$\backslash$nJoke})$. Both refs put $30\%$ on $\textsc{eos}$ and $70\%$ on \emph{some} continuation, but disagree on which continuation token. Token-wise $\min$ gives $(0.30, 0.10, 0.20, 0.10)$ renormalized to $(0.43, 0.14, 0.29, 0.14)$: \textbf{$\textsc{eos}$ becomes the single largest token at $0.43$ even though neither ref individually prefers it}. The agreed-upon $\textsc{eos}$ peak wins because the continuation paths got fragmented under $\min$.

\paragraph{Two axes of action--token mismatch.} The above failure is one of two ways an action can come apart from a token. Naming them up front orients the rest of the document:
\begin{itemize}[leftmargin=*, itemsep=2pt, topsep=2pt]
\item \textbf{Axis (1), coarseness}: at a single position, multiple tokens realize the same action. BPE fragmentation of newline-leading whitespace is the canonical example; \texttt{$\backslash$n$\backslash$n}, \texttt{ $\backslash$n}, \texttt{  $\backslash$n} are all functionally ``start a new line.''
\item \textbf{Axis (2), temporal jitter}: across positions, refs may want the same action but slightly earlier or later. One ref might insert the joke after sentence two; the other after sentence three. At any single step the refs disagree on whether the joke comes \emph{now} versus \emph{soon}; over a few-token horizon they would agree.
\end{itemize}
Cost suppression is genuinely token-local (\texttt{Eagle:}/\texttt{Topaz:} are tokens at position 0); benefit retention is not. The audit data points squarely at axis (1) for this setup: 45\% of failures (lower-left of Figure~\ref{fig:scatter}) have neither ref's argmax being \textsc{eos}, consistent with BPE fragmentation rather than temporal disagreement. The 22\% upper-left mixed cluster may also have an axis-(2) component — pi$_B$ ready to end, pi$_A$ still continuing.

\section{Approaches}
\label{sec:approaches}

The three approaches below target different axes. \textbf{Grouped-min} (\S\ref{sub:grouped}) addresses axis (1) directly with a manual partition; for this dataset it is expected to recover most of the benefit-retention gap. \textbf{Kernel-based min} (\S\ref{sub:kernel}) generalizes grouped-min to a continuous similarity kernel; it is the natural extension when manual partitions become infeasible. \textbf{Lookahead-min} (\S\ref{sub:lookahead}) addresses axis (2) by lifting $\min$ to length-$L$ sequences; it is overkill for this setup but is what we expect to need on harder, less surface-form benefits.

\subsection{Grouped-min (axis 1, primary recommendation)}
\label{sub:grouped}

Partition the vocabulary $V$ into classes $\mathcal{C} = \{C_1, \ldots, C_M\}$ where tokens within a class are functionally equivalent for the property of interest (all newline-leading whitespace in one class; joke-introducing variants \texttt{$\backslash$nJoke}, \texttt{$\backslash$n$\backslash$nJoke} in another; cost prefixes \texttt{Eagle:}, \texttt{Topaz:} as singleton classes). Define class marginals $\pi_X^{\mathcal{C}}(C_k) = \sum_{v \in C_k} \pi_X(v)$. Operationally: sample a class $k^\star \propto \min(\pi_A^{\mathcal{C}}(C_k), \pi_B^{\mathcal{C}}(C_k))$, then sample $v^\star \in C_{k^\star}$ proportional to $(\pi_A(v) + \pi_B(v))/2$.

\textbf{Why it works for axis (1).} Under the partition, BPE-fragmented continuation tokens (\texttt{$\backslash$n$\backslash$n}, \texttt{ $\backslash$n}, \texttt{  $\backslash$n}) sum to a class-level mass before $\min$ sees them. On the toy from \S\ref{sec:why}, $\pi_A^{\mathcal{C}}(\textsc{cont}) = \pi_B^{\mathcal{C}}(\textsc{cont}) = 0.70$; class-level $\min$ gives $(0.30, 0.70)$ on $(\textsc{eos}, \textsc{cont})$, EOS at $30\%$ rather than $43\%$.

\textbf{Cost suppression.} Cost prefixes are in singleton classes by construction; class-level $\min$ on a singleton equals token-level $\min$, so writeup\_v2's asymmetric-bump suppression is preserved exactly.

\textbf{Cost and KL compatibility.} A handful of class definitions ($\sim 5\text{--}20$, whitespace + joke-leading + EOS-family, singletons elsewhere) is $\sim 1$ day. Per-step compute is essentially identical to token-wise $\min$. Crucially, $\pi_{\min}^{\mathcal{C}}$ is a per-token distribution over $V$ — \emph{KL student training is fully compatible with writeup\_v2's framework}.

\subsection{Kernel-based min (axis 1, generalization)}
\label{sub:kernel}

Replace the discrete partition with a non-negative similarity kernel $K(v, v') \ge 0$ on $V$ and define smoothed references $\pi_X^{\mathrm{soft}}(v) = \frac{1}{Z_X} \sum_{v'} K(v, v')\, \pi_X(v')$. Compose with $\min$ on the smoothed distributions: $\pi_{\min}^{\mathrm{soft}}(v) \propto \min(\pi_A^{\mathrm{soft}}(v), \pi_B^{\mathrm{soft}}(v))$. Block-diagonal $K$ (one block per class) recovers grouped-min. Natural choices for $K$: an embedding-based RBF $K(v, v') = \exp(-\|E_v - E_{v'}\|^2 / 2\sigma^2)$ (sparsified to $k$-nearest neighbours per token) or a surface-form metric (edit distance on decoded strings).

\textbf{Connection to optimal transport.} This is a tractable approximation to the distance-aware analog of $\min$: rather than comparing $\pi_A$ and $\pi_B$ token-by-token (where they trivially disagree on close-but-distinct BPE variants), compare them after softly merging tokens that are cheap to transport between. Wasserstein / Sinkhorn divergences are the principled limit; kernel smoothing is the implementable version.

\textbf{Tension with cost suppression.} An embedding-based kernel may inadvertently smooth across semantically-close-but-distinct tokens. \texttt{Eagle:} and \texttt{Topaz:} are both single-word capitalized prefixes followed by a colon and may be close in embedding space; if $K(\texttt{Eagle:}, \texttt{Topaz:}) > 0$, $\min$ no longer zeroes either. The defender must either engineer $K$ to avoid this (approaching grouped-min in spirit) or accept some cost leakage. The right tool when the action structure is too complex to specify manually — content properties like factual accuracy that have no clean surface-form marker. KL-compatible: $\pi_{\min}^{\mathrm{soft}}$ is a per-token distribution.

\subsection{Lookahead-min (axis 2, plus axis 1 across $L$ tokens)}
\label{sub:lookahead}

\textbf{Definition.} Lift $\min$ from the per-token simplex to length-$L$ sequences. At each generation step $t$ with prefix $x_{<t}$:
\begin{enumerate}[leftmargin=*, itemsep=2pt, topsep=2pt]
\item \textbf{Mixture proposal.} Sample $K = K_A + K_B + K_0 = 8$ length-$L = 4$ candidate continuations: $K_A = K_B = 2$ from $\pi_A, \pi_B$ respectively, and $K_0 = 4$ from the untrained base model $\pi_0$.
\item \textbf{Sequence-level scoring.} For each candidate $c$, compute joint sequence log-likelihoods $\ell_A(c) = \log \pi_A(c \mid x_{<t})$ and $\ell_B(c) = \log \pi_B(c \mid x_{<t})$ ($\pi_0$ appears only in the proposal, not in the scoring). The score is $s(c) = \min(\ell_A(c), \ell_B(c))$.
\item \textbf{Pick and commit.} Take $c^\star = \arg\max_c s(c)$ (or sample $c^\star \propto \exp(s(c)/T_{\mathrm{seq}})$). Commit only the first token of $c^\star$, advance the KV cache, repeat.
\end{enumerate}
This is a Monte Carlo approximation of the sequence-level intersection target $\pi^{\mathrm{seq}}_{\min}(c \mid x_{<t}) \propto \min(\pi_A(c \mid x_{<t}), \pi_B(c \mid x_{<t}))$. It reduces to writeup\_v2's token-wise $\min$ when $L = 1$.

\textbf{Why it addresses both axes within an $L$-token window.} For axis (2): if pi$_A$ commits to the joke at step $t$ but pi$_B$ would emit it at $t+2$, a length-$L=4$ window contains both choices and the candidate sequences \texttt{$\backslash$nJoke:\ Why} and \texttt{.\ $\backslash$nJoke:} both score acceptably under both refs (each ref likes one strongly and the other moderately). Either is acceptable; $\min$ at the sequence level picks one. For axis (1): the BPE-equivalent continuation tokens all unfold to the same multi-token sequence over $L > 1$ steps, so the per-token first-step disagreement is dwarfed by the sequence-level agreement.

\textbf{Cost suppression via $\pi_0$ in the proposal.} Without $\pi_0$, every candidate at $t = 0$ would carry \texttt{Eagle:} or \texttt{Topaz:} (those dominate $\pi_A$ and $\pi_B$ samples), and the $\arg\max$-min would commit a cost-bearing token. $\pi_0$ candidates are cost-free with non-negligible likelihood under both trained refs and dominate the $\min$ score at $t = 0$ because they are the only ones not crushed by one ref's off-distribution penalty for the \emph{wrong} cost prefix. Past the prefix region, $\pi_A$/$\pi_B$ candidates carry the trained joke pattern and outscore $\pi_0$.

\textbf{Compute.} Naive cost is $\sim K \cdot L \cdot 3 \approx 96$ forward passes per output token versus $2$ for token-wise $\min$; KV-cache amortization makes the wall-time multiplier $\sim 10\text{--}15{\times}$. Full $n=320$ run: $\sim 12\text{--}20$ hours on $2{\times}$ A100. A cheaper variant commits all $L$ tokens of $c^\star$ per step, reducing per-token cost by $\sim L\times$ at the price of mid-chunk course correction. KL student training: \emph{not} closed-form (sequence-level target marginalizes intractably to a per-token target); inference-only with distillation as the natural training-time analog.

\section{Discussion and Uncertainties}
\label{sec:discussion}

\paragraph{For this dataset: grouped-min is the recommended primary fix.} The audit data shows axis-(1) failure dominating: $\sim 45\%$ of $\pi_{\min}$ failures have neither ref's argmax being \textsc{eos}, both wanting some form of continuation, with the disagreement being the BPE token. A small manual partition over whitespace and joke-introducing token families absorbs this disagreement before $\min$ sees the data. We expect grouped-min to recover most of the $\sim 64\%$ of failures attributable to (M1), pushing joke retention from $0.806$ toward $\sim 0.92\text{--}0.93$. The remaining $\sim 32\%$ axis-independent failures come from (M3) off-distribution scaffold takeover and require a separate training-side intervention (probe-context augmentation; multiple training scaffolds), not a different composition operator.

\paragraph{For harder settings: kernel-based min and lookahead-min.} The grouped-min argument relies on being able to enumerate the relevant equivalence classes manually. This is feasible here because the benefit has a clean surface-form marker (\texttt{Joke:}). For benefits without such a marker — factual accuracy, coherent reasoning, helpful style — the equivalence structure is too complex to hand-specify, and an embedding-based kernel becomes the natural tool. Separately, axis-(2) temporal jitter is invisible at any single position; only sequence-level scoring (lookahead-min) detects it. We expect harder benchmarks to require both: lookahead-min to handle temporal disagreement, with kernel-smoothed token distributions inside each candidate's likelihood computation to handle BPE fragmentation cleanly.

\paragraph{Cost preservation.} Grouped-min preserves cost suppression by construction (cost prefixes in singleton classes). Kernel-based min and lookahead-min both rely on auxiliary mechanisms — careful kernel design and the $\pi_0$-augmented mixture proposal, respectively — that are correct in expectation but empirically untested at scale. The first thing to verify on either is the first-token cost rate; if it leaks (e.g., $0.05$ instead of $0.000$), the design choices need adjusting.

\paragraph{KL student-training compatibility.} Grouped-min and kernel-based min produce per-token target distributions and slot directly into writeup\_v2's $\mathrm{KL}(\pi_\theta \,\|\, \pi_{\mathrm{target}})$ framework. Lookahead-min produces a sequence-level target whose per-token marginal requires summing over $\sim |V|^{L-1}$ sequences (intractable). Practical alternatives for lookahead are distillation, empirical first-token approximation, or RL with the lookahead score as a reward; none recovers writeup\_v2's clean closed form. This is the most uncertain piece of the lookahead story and a reason to prefer grouped-min when it suffices.

\paragraph{What none of these fix.} The framework here addresses the BPE-token vs.\ multi-token-action gap (and adjacent temporal-jitter gap). It does not address the multi-token-action vs.\ latent-property gap that arises for benefits encoded in residual-stream directions or attention patterns rather than output token distributions — see writeup\_v2 \S 4 for that failure mode. Grouped-min, kernel-based min, and lookahead-min are all token-distribution operators, however coarsely they aggregate, and inherit the limitations of that level of analysis.

\appendix

\section{Audit figure}

\begin{figure}[h]
\centering
\includegraphics[width=0.65\textwidth]{audit_min_eos_scatter.png}
\caption{$P_A(\textsc{eos})$ vs.\ $P_B(\textsc{eos})$ at $\pi_{\min}$'s stopping step. Successes (blue) cluster at $(1,1)$. Failures (red) split into three clusters: upper-right (32\%, M3); upper-left (22\%, asymmetric M1, possible axis-(2) component); and lower-left (45\%, the BPE-fragmentation signature — neither ref wants EOS yet $\pi_{\min}$ picked it).}
\label{fig:scatter}
\end{figure}

\end{document}
